\begin{definition} \hlt{Block Matrices}\\
Suppose $\mathcal{V} = \mathcal{U}_1 \oplus \mathcal{U}_2$ and $\mathcal{W} = \mathcal{Z}_1 \oplus \mathcal{Z}_2$. If $T \in \LL(\mathcal{V}, \mathcal{W})$, relative to bases adapted to these decompositions in the sense of Definition \ref{def:adapted-basis-direct-sum}, the matrix of $T$ has block form $[T] =
\begin{bmatrix}
A_{11} & A_{12} \\
A_{21} & A_{22}
\end{bmatrix}$,
where each block records the part of $T$ mapping one summand of the domain into one summand of the codomain.
\end{definition}

\begin{definition} \hlt{Conformal Block Partitions}\\
For a block matrix $A=[A_{ij}]$, the submatrices $A_{ij}$ are called \hlt{blocks}. Two block partitions are \hlt{conformal for addition} when corresponding blocks have the same sizes. They are \hlt{conformal for multiplication} when every block product appearing in the usual row-column block multiplication is defined.
\end{definition}

\begin{proposition} \label{prop:block-matrix-arithmetic} \hlt{Block Matrix Arithmetic}\\
For conformal $2\times2$ block matrices,
\begin{align}
\begin{bmatrix}
A_{11} & A_{12}\\
A_{21} & A_{22}
\end{bmatrix}
+
\begin{bmatrix}
B_{11} & B_{12}\\
B_{21} & B_{22}
\end{bmatrix}
=
\begin{bmatrix}
A_{11}+B_{11} & A_{12}+B_{12}\\
A_{21}+B_{21} & A_{22}+B_{22}
\end{bmatrix}
\nonumber
\end{align}
and
\begin{align}
\begin{bmatrix}
A_{11} & A_{12}\\
A_{21} & A_{22}
\end{bmatrix}
\begin{bmatrix}
B_{11} & B_{12}\\
B_{21} & B_{22}
\end{bmatrix}
=
\begin{bmatrix}
A_{11}B_{11}+A_{12}B_{21} & A_{11}B_{12}+A_{12}B_{22}\\
A_{21}B_{11}+A_{22}B_{21} & A_{21}B_{12}+A_{22}B_{22}
\end{bmatrix}.
\nonumber
\end{align}
\begin{proof}
Each displayed block consists of the same entries obtained by ordinary entrywise addition or ordinary row-column multiplication, with the common indices grouped according to the partition.\\
\hlt{Structural Commentary}: This is ordinary matrix arithmetic with indices grouped into rectangular packages.\\
\hlt{Key Insight}: Conformality is exactly the condition that the grouped row-column products make sense.
\end{proof}
\end{proposition}

\begin{proposition} \label{prop:block-centrosymmetric-commutation} \hlt{Block Centrosymmetry and Commutation}\\
Let $M=
\begin{bmatrix}
A & B\\
C & D
\end{bmatrix}\in\operatorname{Mat}(2n,2n,\mathbb{F})$,
$L_{2n}=
\begin{bmatrix}
0 & I_n\\
I_n & 0
\end{bmatrix}$.\\
Then $M$ has block centrosymmetric form $A=D$ and $B=C$ if and only if $L_{2n}M=ML_{2n}$.
\begin{proof}
Block multiplication gives $L_{2n}M=
\begin{bmatrix}
C & D\\
A & B
\end{bmatrix}$ and 
$ML_{2n}=
\begin{bmatrix}
B & A\\
D & C
\end{bmatrix}$.\\
The two matrices are equal exactly when $A=D$ and $B=C$.\\
\hlt{Structural Commentary}: The matrix $L_{2n}$ swaps the two block coordinates. Commuting with that swap means $M$ treats the two block positions symmetrically.\\
\hlt{Key Insight}: Compute both block products; equality forces matching opposite blocks.
\end{proof}
\end{proposition}

\begin{definition} \hlt{Direct Sum of Matrices}\\
For $A\in\operatorname{Mat}(m,n,\mathbb{F})$ and $B\in\operatorname{Mat}(p,q,\mathbb{F})$, the \hlt{direct sum} of $A$ and $B$ is $A\oplus B\defeq
\begin{bmatrix}
A & 0\\
0 & B
\end{bmatrix}
\in\operatorname{Mat}(m+p,n+q,\mathbb{F})$.
\end{definition}

\begin{definition} \label{def:polynomial-evaluation-on-square-matrix} \hlt{Polynomial Evaluation on a Square Matrix}\\
If $A\in\operatorname{Mat}(n,n,\mathbb{F})$ and $p(z)=\sum\limits_{j=0}^d c_jz^j\in\mathbb{F}[z]$, then
\[
p(A)\defeq\sum\limits_{j=0}^d c_jA^j,
\qquad A^0\defeq I_n.
\]
\end{definition}

\begin{proposition} \label{prop:direct-sum-matrix-algebra} \hlt{Algebra of Matrix Direct Sums}\\
Whenever the expressions are defined,
\begin{enumerate}[label=(\roman*)]
\setlength{\itemsep}{0pt}
\item $a(A\oplus B)+b(C\oplus D)=(aA+bC)\oplus(aB+bD)$;
\item $(A\oplus B)(C\oplus D)=AC\oplus BD$;
\item if $A$ and $B$ are square, then $(A\oplus B)^k=A^k\oplus B^k$ for every positive integer $k$;
\item if $A$ and $B$ are square and $p$ is a polynomial, then $p(A\oplus B)=p(A)\oplus p(B)$.
\end{enumerate}
\begin{proof}
Block addition gives $a(A\oplus B)+b(C\oplus D)
=
\begin{bmatrix}
aA+bC&0\\
0&aB+bD
\end{bmatrix}
=(aA+bC)\oplus(aB+bD)$.\\
Block multiplication gives $(A\oplus B)(C\oplus D)
=
\begin{bmatrix}
AC&0\\
0&BD
\end{bmatrix}
=AC\oplus BD$.\\
For powers, the case $k=1$ is the definition. If $(A\oplus B)^k=A^k\oplus B^k$, then part (ii) gives $(A\oplus B)^{k+1}=A^{k+1}\oplus B^{k+1}$. Hence (iii) follows by induction. For the polynomial statement, suppose $A\in\operatorname{Mat}(m,m,\mathbb{F})$ and $B\in\operatorname{Mat}(r,r,\mathbb{F})$. For $j=0$, Definition \ref{def:polynomial-evaluation-on-square-matrix} gives $(A\oplus B)^0=I_{m+r}=I_m\oplus I_r=A^0\oplus B^0$. If $p(z)=\sum_{j=0}^{d}c_jz^j$,\\
then $p(A\oplus B)=\sum\limits_{j=0}^{d}c_j(A^j\oplus B^j) =\left(\sum\limits_{j=0}^{d}c_jA^j\right)\oplus\left(\sum\limits_{j=0}^{d}c_jB^j\right) =p(A)\oplus p(B)$.\\
\hlt{Structural Commentary}: A direct sum matrix keeps the two blocks independent under linear combinations, products, powers, and polynomials.\\
\hlt{Key Insight}: Block diagonal matrices behave componentwise.
\end{proof}
\end{proposition}

\begin{proposition} \label{prop:diagonal-multiplication-scales-rows-columns} \hlt{Diagonal Multiplication Scales Rows and Columns}\\
Let $A=[a_{ij}]\in\operatorname{Mat}(m,n,\mathbb{F})$, let $\Lambda=\operatorname{diag}(\lambda_1,\ldots,\lambda_n)$, and let $M=\operatorname{diag}(\mu_1,\ldots,\mu_m)$. Then $A\Lambda=[\lambda_j a_{ij}]$ and $MA=[\mu_i a_{ij}]$. Thus right multiplication by $\Lambda$ scales columns of $A$, left multiplication by $M$ scales rows of $A$.
\begin{proof}
The $(i,j)$ entry of $A\Lambda$ is $\sum_{r=1}^n a_{ir}\Lambda_{rj}=a_{ij}\lambda_j$, and the $(i,j)$ entry of $MA$ is $\sum_{r=1}^m M_{ir}a_{rj}=\mu_i a_{ij}$.\\
\hlt{Structural Commentary}: Diagonal matrices act by independent scalar weights on coordinate directions.\\
\hlt{Key Insight}: Multiplication on the right scales columns, while multiplication on the left scales rows.
\end{proof}
\end{proposition}

\begin{proposition} \label{prop:block-upper-triangular-inverse} \hlt{Inverse of an Upper Block-Triangular Matrix}\\
Let $Y\in\operatorname{Mat}(m,m,\mathbb{F})$ and $Z\in\operatorname{Mat}(n,n,\mathbb{F})$ be invertible, and let $X\in\operatorname{Mat}(m,n,\mathbb{F})$. Then $\begin{bmatrix}
Y & X\\
0 & Z
\end{bmatrix}^{-1}
=
\begin{bmatrix}
Y^{-1} & -Y^{-1}XZ^{-1}\\
0 & Z^{-1}
\end{bmatrix}$.
\begin{proof}
Let $M=\begin{bmatrix}Y&X\\0&Z\end{bmatrix}$ and $N=\begin{bmatrix}Y^{-1}&-Y^{-1}XZ^{-1}\\0&Z^{-1}\end{bmatrix}$. Proposition \ref{prop:block-matrix-arithmetic} gives
\begin{align}
MN=
\begin{bmatrix}
I_m&-XZ^{-1}+XZ^{-1}\\
0&I_n
\end{bmatrix}
=I_{m+n},
\qquad
NM=
\begin{bmatrix}
I_m&Y^{-1}X-Y^{-1}X\\
0&I_n
\end{bmatrix}
=I_{m+n} \nonumber
\end{align}
Thus $N=M^{-1}$.\\
\hlt{Structural Commentary}: Upper block-triangular inversion preserves the zero lower-left block. The off-diagonal term is forced by cancelling the original upper-right block.\\
\hlt{Key Insight}: Inverting a block triangular matrix inverts diagonal blocks and corrects the coupling block.
\end{proof}
\end{proposition}

\begin{theorem} \label{thm:inverse-upper-triangular} \hlt{Inverse of an Upper Triangular Matrix}\\
Let $A=[a_{ij}]\in\operatorname{Mat}(n,n,\mathbb{F})$ be upper triangular with nonzero diagonal entries. Then $A$ is invertible, $A^{-1}$ is upper triangular, and the diagonal entries of $A^{-1}$ are $a_{11}^{-1},\ldots,a_{nn}^{-1}$.
\begin{proof}
We argue by induction on $n$. The case $n=1$ is immediate. For $n>1$, write $A=
\begin{bmatrix}
B & \mathbf{x}\\
0 & a_{nn}
\end{bmatrix}$, where $B$ is $(n-1)\times(n-1)$ upper triangular with nonzero diagonal entries.\\
By induction, $B^{-1}$ is upper triangular with diagonal entries $a_{11}^{-1},\ldots,a_{n-1,n-1}^{-1}$.\\
Proposition \ref{prop:block-upper-triangular-inverse} gives $A^{-1}=
\begin{bmatrix}
B^{-1} & -B^{-1}\mathbf{x}a_{nn}^{-1}\\
0 & a_{nn}^{-1}
\end{bmatrix}$,
which is upper triangular and has the asserted diagonal.\\
\hlt{Structural Commentary}: An upper triangular matrix splits into a smaller upper triangular block and one final scalar block. The block inverse formula preserves this triangular structure.\\
\hlt{Key Insight}: Nonzero diagonals make blockwise triangular inversion possible.
\end{proof}
\end{theorem}

\begin{proposition} \label{prop:upper-block-triangular-similarity-shift} \hlt{Upper Block-Triangular Similarity Shift}\\
Let $B\in\operatorname{Mat}(m,m,\mathbb{F})$, let $D\in\operatorname{Mat}(n,n,\mathbb{F})$, and let $C,X\in\operatorname{Mat}(m,n,\mathbb{F})$.\\
Then $\begin{bmatrix}
B & C\\
0 & D
\end{bmatrix}
$ is similar to $\begin{bmatrix}
B & C+XD-BX\\
0 & D
\end{bmatrix}$.
\begin{proof}
The matrices $\begin{bmatrix} I_m & X\\ 0 & I_n \end{bmatrix}$ and $\begin{bmatrix} I_m & -X\\ 0 & I_n \end{bmatrix}$ are inverses in $\operatorname{Mat}(m+n,m+n,\mathbb{F})$, and block multiplication gives
\begin{align}
\begin{bmatrix}
I_m & X\\
0 & I_n
\end{bmatrix}
\begin{bmatrix}
B & C\\
0 & D
\end{bmatrix}
\begin{bmatrix}
I_m & -X\\
0 & I_n
\end{bmatrix}
=
\begin{bmatrix}
B & C+XD-BX\\
0 & D
\end{bmatrix}.
\nonumber
\end{align}
\hlt{Structural Commentary}: Unit upper block-triangular conjugation changes only the coupling block.\\
\hlt{Key Insight}: Similarity can adjust the connecting block while preserving the diagonal blocks.
\end{proof}
\end{proposition}
